% Source: https://artofproblemsolving.com/community/c3h2514092_an_interesting_way_of_generating_pythagorean_triples
% Stable identity: aops-topic-2514092-q01
\textbf{(a)} The triple $(9,12,15)$ is a counterexample: $81+144=225$. In the claimed form one leg is $2ab$, which is even, so $2ab=12$ and $ab=6$. With $a>b$ this gives $(a,b)=(6,1)$ or $(3,2)$, which produce the triples $(35,12,37)$ and $(5,12,13)$, not $(9,12,15)$. So $(9,12,15)$ is not of the form. (Every triple is a multiple of a triple of the form, e.g. $(9,12,15)=3\cdot(3,4,5)$ with $(a,b)=(2,1)$, but the multiples are missed.)

\textbf{(b)} Take $a=3$, $b=1$, which are coprime: the triple is $(8,6,10)$, and all three numbers are even. So $\gcd(a,b)=1$ is not enough.

Now let $\gcd(a,b)=1$ with $a$, $b$ of opposite parity, and suppose a prime $p$ divides all three numbers. Then $p$ divides $\left(a^2+b^2\right)+\left(a^2-b^2\right)=2a^2$ and $\left(a^2+b^2\right)-\left(a^2-b^2\right)=2b^2$. Since $a$ and $b$ have opposite parity, $a^2-b^2$ is odd, so $p\ne2$. Hence $p\mid a^2$ and $p\mid b^2$, so $p\mid a$ and $p\mid b$, contradicting $\gcd(a,b)=1$. So the greatest common divisor of the triple is $1$ and the triple is primitive. $\blacksquare$

(The parity condition is needed: if $a$ and $b$ are both odd, all three numbers are even, as in the counterexample.)
